% page 534 of the facsimile (image p-533 of the scan)
% block 154.9 x 242.0 mm ; left margin 8.0 mm
\pgc{-12.700mm}{0.000mm}{
\l{\cel{3}526.}
\l{}
\l{\textquotedbl{}\sou{si\textquotedbl{}.}\cel{1}- (muziko).La noto sepesma di la \textquotedbl{}do\textquotedbl{}gamo..}
\l{}
\l{+si\cel{4}Adverbo per qua onu afirmas la kontreajo di lo juss}
\l{\cel{4}dicita. (ex.: \textquotedbl{}Me esas +certena ke vu ne venos.\textquotedbl{} - \textquotedbl{}Si,}
\l{\cel{3}me venos.\textquotedbl{} - \textquotedbl{}No !\textquotedbl{} - \textquotedbl{}Si !\textquotedbl{}). - F..}
\l{}
\l{\cel{1}sibarito. Persono qua esforcas obtenar komforto maxima por}
\l{\cel{3}su ipsa, e di qua la sentemeso korpala esas rafinita.--}
\l{\cel{4}DEFIRS.}
\l{}
\l{sibilo. (en\cel{1}\sou{la epoki antiqua)}. Muliero pri qua onu asertiss}
\l{\cel{4}ke lu recevas de deo la talento dicar to quo esas even-}
\l{\cel{4}tonta. - DEFIRS.}
\l{}
\l{+\sou{sij)o}. Loko ube ulu, ulo, esas establisita- ube eventas lua}
\l{\cel{4}operaci komercala, industriala, financala. - F.}
\l{}
\l{siena-tero. Okro bruna, redatra, uzata por facar ula farbi..}
\l{}
\l{\cel{1}\sou{siesaar}. (\sou{netrans.}) Dormar pos la repasto dimezala, dejuna.}
\l{\cel{4}- DEFIS.}
\l{}
\l{\cel{1}\sou{sifiliso}. (\sou{patol.}) Morbo infektiva e kontagiala, transmiseba}
\l{\cel{4}a la decedonti, e di qua la kauzo esas la mikrobo L.\cel{1}\sou{tre}-}
\l{\cel{4}\sou{ponema pallidum}\cel{1}- morbo qua afektas precipue la organi}
\l{\cel{3}sexuala. - DEFIRS..}
\l{}
\l{siflar. (trans. e netrans.) Produktar bruiso akuta per igarr}
\l{\cel{4}aero eskapar per aperturo streta. - F. (Kuglo\cel{1}\sou{sisas)}.}
\l{}
\l{sifono.\cel{2}(tekn.)\cel{2}Aparato U-forma, kun la U-branchi longaa}
\l{\cel{4}ne-inter-egale, uzata por transvarsar la liquido di vazo}
\l{\cel{4}aden qua esas imersita la extremajo di la brancho kurta}
\l{\cel{3}aden vazo ube abutas la brancho longa - la atmosferepresoo}
\l{\cel{4}esante plu granda sur la facio ube la kolono de liquido}
\l{\cel{4}esas min alta kande la tubo esas amorcita. - DEFIRS.}
\l{}
\l{\cel{1}\sou{sigaro}. Mikra rulo-bloko ek tabako-folii, di qua onu acendas}
\l{\cel{4}un ek la extremaji, e de qua onu aspiras la fumuro per}
\l{\cel{4}la altra extremajo quan onu pozas en sua boko.- DEFIRS.}
\l{}
\l{sigilario. (bot.) Genero de kriptogamo vaskuloza, sorto dii}
\l{\cel{4}likopodo giganta, arboro ek la plantaro karbon-mineyala,}
\l{\cel{4}di qua la staturo povas atingar 30 metri. - DEFIRS.}
\l{}
\l{siglar.\cel{2}(trans.)\cel{2}Klozar per vaxo markizita (kande lu esiss}
\l{\cel{4}ankore mola pro fuzeso) konkave o reliefe, per apliko}
\l{\cel{4}di korpo harda. - DIS.}
\l{}
\l{\cel{1}\sou{signo}. I. Kozo perceptebla qua venigas la ideo di ento o di}
\l{\cel{3}eso-maniero, konseque de interrelato naturala o konven--}
\l{\cel{4}cionala. II. La\cel{1}\sou{signi zodiakala}\cel{1}: la dek-e-du stelari qui}
\l{\cel{3}formacas zodiako. - III.Gesto per qua onu manifestas too}
\l{\cel{4}quon onu sentas, pensas, od opinionas. - EFIS.}
}
